% ============================================================================
%  CHAPTER 6: TOOLS, TECHNOLOGIES, AND PARAMETERS   (expanded version)
%  Replaces the entire Chapter 6 body in main.tex.
% ============================================================================

\chapter{Tools, Technologies, and Parameters}
\label{chap:tools}

FinITR-AI is written entirely in Python and is built to run on a single ordinary
computer with no cloud dependency, and that constraint --- offline operation on
commodity hardware --- shaped nearly every technology choice in the project.
Wherever a task called for a language model, a local model served through Ollama
was used in place of a hosted API. Wherever embeddings or classification were
needed, the library chosen was one that runs comfortably on a CPU. And wherever
arithmetic was involved, it was written from scratch rather than delegated to any
external service. This chapter records the concrete stack, with the reasoning
behind the choices that were not obvious, and then sets out the numeric
parameters that govern the system's behaviour so that the results in the
previous chapter are reproducible.

\section{Tools and Technologies}
\label{sec:tools}

\begin{itemize}
    \item \textbf{Programming language and application frameworks.} Python 3.10
          or later throughout. The backend is a FastAPI service run under
          Uvicorn, exposing endpoints for analysis, the interview step, and the
          report exports; the interactive front-end is built in Streamlit; and
          Pydantic provides the request and response schemas that keep the two
          ends in agreement.
    \item \textbf{Hardware and platform.} The system targets commodity hardware
          --- a minimum of 8\,GB of RAM, with 16\,GB recommended. A GPU is
          optional; an RTX 3060 or better speeds up local model inference but is
          not required for correctness, only for latency. Everything is designed
          to run offline on the user's own device, which is the property that
          makes it safe for sensitive documents.
    \item \textbf{Local large language model.} A model served through Ollama,
          with qwen2.5:7b as the recommended default and llama3.1:8b as a tested
          alternative. The model is used only for natural-language explanation
          and for the rare classification fallback; it never performs
          arithmetic, and the pipeline degrades to deterministic templates when
          it is unavailable, so no result depends on the model being up.
    \item \textbf{Machine-learning and NLP libraries.} The transaction classifier
          embeds descriptions with a multilingual MiniLM sentence-transformer
          \cite{ref10}, chosen because it handles Hinglish merchant names on a
          CPU; ChromaDB serves as the vector store for semantic document search;
          scikit-learn \cite{ref11} provides both the logistic-regression notice
          predictor and the k-nearest-neighbour stage of the classifier; and the
          faithfulness verifier wraps a DeBERTa-based NLI cross-encoder
          \cite{ref9}. NumPy underlies the numerical work.
    \item \textbf{Document parsing.} pandas handles the bank-statement CSV and
          other tabular data, while pdfplumber and pdfminer.six extract text from
          Form 16 PDFs. A dedicated description normaliser strips transfer codes,
          UTR references, and transaction IDs before classification.
    \item \textbf{Output generation and visualisation.} ReportLab and Jinja2
          generate the PDF CA-brief; plotly renders the dashboard visualisations
          --- the risk gauge, the reconciliation table, and the regime Sankey
          diagram --- in the Streamlit front-end.
    \item \textbf{Datasets and external resources.} IndianTaxBench v1.0, the
          purpose-built suite of 140 labelled synthetic cases (100 development
          and 40 held-out) across 22 employers and 100 PANs, together with the
          synthetic Form 16, AIS, and bank-statement fixtures used in
          development; and a hand-built Income-Tax-Act knowledge corpus of more
          than eighty section nodes that the PageIndex retriever walks. The
          feature design for the notice predictor draws on the CBDT Annual Report
          \cite{ref12}, which identifies AIS mismatch as the leading cause of
          Section 143(1) adjustments.
    \item \textbf{Development and testing tools.} pytest for unit and integration
          tests, httpx for exercising the API, tqdm for evaluation progress, and
          standard Python virtual environments for dependency isolation.
\end{itemize}

\section{Key Parameters and Settings}
\label{sec:parameters}

Table~\ref{tab:parameters} lists the parameters that most directly determine the
system's behaviour and its measured results. They fall into two groups. Some are
statutory values for FY 2025-26 --- these are not tunable, and getting them
exactly right is the whole point of the deterministic engine. The rest are
genuine tuning choices, and almost all of them reflect the false-negative-averse
philosophy described in Chapter~\ref{chap:challenges}: where a parameter trades
one kind of error against another, it is set to prefer the cheaper mistake.

\begin{table}[H]
\centering
\caption{Important Parameters Used in the Implementation}
\label{tab:parameters}
\small
\begin{tabularx}{\textwidth}{|l|X|}
\hline
\textbf{Parameter} & \textbf{Value / Rationale} \\
\hline
Orchestrator max iterations & 3 --- bounds the critic re-run loop; in practice cases resolve in one or two passes, and the bound prevents runaway self-correction \\
\hline
Form-16 vs AIS noise tolerance & 3\% --- below this gap the AIS salary is trusted as ground truth, absorbing OCR and rounding noise \\
\hline
New Regime standard deduction & Rs.~75,000 (FY 2025-26, Section 115BAC) --- statutory \\
\hline
Old Regime standard deduction & Rs.~50,000 (FY 2025-26) --- statutory \\
\hline
Health \& education cess & 4\% on tax plus surcharge, both regimes --- statutory \\
\hline
Section 87A rebate (new) & up to Rs.~60,000 for taxable income $\leq$ Rs.~12,00,000, with marginal relief to $\approx$ Rs.~12,75,000 --- statutory \\
\hline
Section 87A rebate (old) & up to Rs.~12,500 for taxable income $\leq$ Rs.~5,00,000 --- statutory \\
\hline
Risk weights (additive) & crypto = 60, capital gains = 55, AIS mismatch = 25, freelance = 25, property = 20, cash = 15, savings interest = 10 \\
\hline
Risk bands & LOW $<$20, MEDIUM $<$50, HIGH $<$75, CRITICAL $\geq$75 (score clipped to 100) \\
\hline
Classifier confidence threshold & 0.65 --- below this the kNN stage defers to the LLM fallback \\
\hline
Classifier neighbours (k) & 5 --- majority vote over the five nearest labelled descriptions \\
\hline
Notice-predictor threshold & 0.032 --- lowest threshold with notice-recall $\geq$ 0.90; test AUC 0.9524 \\
\hline
Notice-predictor model & Logistic Regression over 14 features (chosen over GBC, which overfits at this scale: test AUC 0.875) \\
\hline
Transaction embedding model & multilingual MiniLM (384-dim), CPU-friendly \\
\hline
NLI verifier model & DeBERTa-v3 cross-encoder; keyword-overlap fallback if unavailable \\
\hline
Default local LLM & qwen2.5:7b via Ollama \\
\hline
\end{tabularx}
\end{table}
